\documentclass[11pt]{article}

\usepackage[T1]{fontenc}
\usepackage[utf8]{inputenc}
\usepackage[spanish,es-nodecimaldot]{babel}
\usepackage{amsmath,amssymb,mathtools}
\usepackage{enumitem}
\usepackage[a4paper,margin=2.5cm]{geometry}
\usepackage[hidelinks]{hyperref}

\newcommand{\Cstar}{C^*}
\newcommand{\Asa}{A_{\mathrm{sa}}}
\newcommand{\U}{\mathcal{U}}
\newcommand{\Pj}{\mathcal{P}}
\newcommand{\one}{\mathbf{1}}

\newenvironment{definicion}[1]
  {\par\medskip\noindent\textbf{Definición #1.}\ }
  {\par\medskip}

\title{Definiciones de la Parte 2: $C^*$-álgebras}
\author{}
\date{}

\begin{document}
\maketitle

\section*{Índice de definiciones}
\begin{enumerate}[label=\textbf{Definición 2.\arabic*.},leftmargin=*,start=1]
  \item $C^*$-álgebra. \hfill (p. 1)
\end{enumerate}
\begin{enumerate}[label=\textbf{Definición 2.\arabic*.},leftmargin=*,start=4]
  \item $*$-homomorfismo. \hfill (p. 1)
  \item Elementos normales, autoadjuntos, unitarios y proyecciones. \hfill (p. 1)
\end{enumerate}
\begin{enumerate}[label=\textbf{Definición 2.\arabic*.},leftmargin=*,start=23]
  \item Extensión del cálculo funcional. \hfill (p. 6)
\end{enumerate}

\medskip
\noindent\textit{Nota.} El documento original pasa de la Definición/Observación 2.2 a la Definición 2.4; se conserva exactamente esa numeración y no se introduce una ``Definición 2.3''.

\section*{Definiciones}

% Fuente: Parte2.pdf, p. 1
\begin{definicion}{2.1}
Una $C^*$-álgebra $A$ es un álgebra de Banach involutiva tal que
\[
  \lVert a^*a\rVert = \lVert a\rVert^2,
  \qquad \forall a\in A,
\]
(identidad $C^*$).
\end{definicion}

% Fuente: Parte2.pdf, p. 1
\begin{definicion}{2.4}
Sean $A$ y $B$ $C^*$-álgebras. Un $*$-homomorfismo es un mapeo
\[
  \varphi:A\longrightarrow B
\]
lineal, multiplicativo y que preserva la involución,
\[
  \varphi(a^*)=\varphi(a)^*.
\]
Si $A$ y $B$ tienen unidad, decimos que $\varphi$ es \emph{unital} si
\[
  \varphi(1_A)=1_B.
\]
\end{definicion}

% Fuente: Parte2.pdf, p. 1
\begin{definicion}{2.5}
Sea $a\in A$. Decimos que $a$ es:
\begin{enumerate}[label=(\roman*),leftmargin=2.2em]
  \item \emph{normal} si
  \[
    a^*a=aa^*;
  \]

  \item \emph{autoadjunto} si
  \[
    a^*=a.
  \]
  Denotamos al conjunto de autoadjuntos con $\Asa$;

  \item \emph{unitario} si
  \[
    a^*a=aa^*=1
  \]
  (asumiendo que $A$ tiene unidad). Escribimos
  \[
    \U(A)=\{\text{unitarios de }A\};
  \]

  \item \emph{proyección} si
  \[
    a=a^*=a^2.
  \]
  Escribimos
  \[
    \Pj(A)=\{\text{proyecciones de }A\}.
  \]
\end{enumerate}
\end{definicion}

% Fuente: Parte2.pdf, p. 6
\begin{definicion}{2.23}
Sea $A$ con unidad y $a\in A$ normal. Sea $E\subseteq\mathbb{C}$ un espacio localmente compacto que contiene a $\sigma(a)$. Definimos un $*$-homomorfismo unital
\[
  C_b(E)\longrightarrow A,
  \qquad
  f\longmapsto \bigl(f|_{\sigma(a)}\bigr)(a)=:f(a),
\]
utilizando el cálculo funcional. Si $f:E\to\mathbb{C}$ no es acotada, del mismo modo se define $f(a)\in A$ (pues lo único que importa es $f|_{\sigma(a)}$, que sí es acotada por ser $\sigma(a)$ compacto).
\end{definicion}

\end{document}
